\section{Irrational slopes for \texorpdfstring{$m\ge4$}{m >= 4}}\label{sec:irrational}
\subsection{Existence of the limit}

\begin{lemma}[Vertical length comparison]\label{lem:vertical-length-comparison}
Let $M\in\Z_{>0}$ and $L\in\Z_{>0}$. Let $X_1,\ldots,X_L$ and $Y_1,\ldots,Y_L$ be simple runs.
Put $X=X_1\ccn \cdots\ccn X_L$, $Y=Y_1\ccn \cdots\ccn Y_L$, $\alpha_i=|X_i|_N$, and $\beta_i=|Y_i|_N$.
Assume $\alpha_i\in[M]_{\Z}$ and $\beta_i\in[M]_{\Z}$ for $i\in[L]_{\Z}$.
Then, there are $n\in\Z_{\ge0}$ and a run-edit list $\Omega=(\omega_1,\ldots,\omega_n)$ satisfying the following.
\begin{enumerate}
\item\label{c:vertical-comparison-target} The list $\Omega$ starts at $X$ and satisfies
\[
 \operatorname{RunEResult}(\Omega,X)
 =\operatorname{Trans}\bigl(Y,\operatorname{Start}_{\mathrm{pt}}(X)-\operatorname{Start}_{\mathrm{pt}}(Y)\bigr).
\]
\item\label{c:vertical-comparison-count} The list $\Omega$ consists of exactly
\[
 n=\sum_{i=1}^L|\alpha_i-\beta_i|
\]
vertical run-edits.
\item\label{c:vertical-comparison-lengths} After each run-edit of $\Omega$, every simple-run vertical length belongs to $[M]_{\Z}$.
\end{enumerate}
\end{lemma}
\begin{proof}
\prf{Claim~\ref{c:vertical-comparison-target}} Let us construct $\Omega$.
Process the simple runs in the order $i=1,\ldots,L$.
At the $i$th simple run, let $W$ be the current path and write this simple run as $W[s,s+t+1]$.
Here $s\in[0,|W|-1]_{\Z}$ and $t\in\Z_{>0}$.

First, let us increase the vertical length when $\alpha_i<\beta_i$.
For each $t\in[\alpha_i,\beta_i-1]_{\Z}$ in increasing order, replace $W$ by
\[
 \operatorname{Ins}(W,s+t+1,N).
\]
This gives $\beta_i-\alpha_i$ insertions and final vertical length $\beta_i$.

Second, let us decrease the vertical length when $\beta_i<\alpha_i$.
For each $t=\alpha_i,\alpha_i-1,\ldots,\beta_i+1$, replace $W$ by
\[
 \operatorname{Del}(W,W_{s+t+1}).
\]
Since $t>\beta_i\ge1$, every deletion has $t\ge2$.
This gives $\alpha_i-\beta_i$ deletions and final vertical length $\beta_i$.

Let us identify the resulting path.
When $\alpha_i=\beta_i$, make no change at the $i$th simple run.
The resulting vertical lengths are $\beta_1,\ldots,\beta_L$.
Thus, the resulting path is $\operatorname{Trans}(Y,\operatorname{Start}_{\mathrm{pt}}(X)-\operatorname{Start}_{\mathrm{pt}}(Y))$.

\prf{Claim~\ref{c:vertical-comparison-count}} Let us count the vertical run-edits.
For each $i\in[L]_{\Z}$, the construction in Claim~\ref{c:vertical-comparison-target} uses exactly $|\alpha_i-\beta_i|$ vertical run-edits. Summing over $i\in[L]_{\Z}$ gives $n=\sum_{i=1}^L|\alpha_i-\beta_i|$.

\prf{Claim~\ref{c:vertical-comparison-lengths}} Let us bound the intermediate vertical lengths.
During the run-edits of $\Omega$ on the $i$th simple run, its vertical length $t$ satisfies
\[
 1\le\min(\alpha_i,\beta_i)\le t\le\max(\alpha_i,\beta_i)\le M.
\]
For $j\in[L]_{\Z}\setminus\{i\}$, these run-edits do not change the vertical length of the $j$th simple run.
\end{proof}

\Needspace{15\baselineskip}
\begin{definition}\label{def:admissible-mechanical}
Let $M\in\Z_{>0}$, $r\in\R$, $\eta>0$, $\delta>0$, and $B>0$.
Let $l\in\Z_{>0}$, $T\in\Z_{>0}$, $\sigma_1\in[1,\infty)$, $\sigma_2\in[1,\infty)$,
$u_1\in[0,1)$, and $u_2\in[0,1)$. Let $A_1,\ldots,A_l$ be simple runs.
We call $(l,T,\sigma_1,\sigma_2,u_1,u_2,(A_1,\ldots,A_l))$
\textit{$(M,r,\eta,\delta,B)$-admissible} when the following conditions hold.
\begin{hypotheses}
\item\label{hyp:mechanical-source} We have
\[
 |A_i|_N=\floor{u_1+i\sigma_1}-\floor{u_1+(i-1)\sigma_1}
 \quad(i\in[l-1]_{\Z}),\qquad |A_l|_N\in[M]_{\Z}.
\]
\item\label{hyp:mechanical-proximity} We have
\[
 |\sigma_i-r|<\delta\quad(i\in[2]_{\Z}),\qquad T\ge(1+\eta)l+B.
\]
\end{hypotheses}
\end{definition}

\begin{lemma}[Fractional parts near a rational slope]\label{lem:fractional-part-estimates}
Let $p\in\Z_{>0}$ and $q\in\Z_{>0}$ satisfy $\gcd(p,q)=1$.
Let $\sigma\in\R$, $w\in\R$, and $0<\eta<\frac14$. Assume
\[
 |\sigma-\frac pq|<\frac2{q^2}.
\]
For $x\in\R$, let $\{x\}=x-\floor x$.
Let $0\le x_1\le\cdots\le x_q<1$ be the nondecreasing rearrangement of
$(\{w+j\sigma\})_{j\in[0,q-1]_{\Z}}$. Let $x_{q+1}=x_1+1$.
Then, the following hold.
\begin{enumerate}
\item\label{c:fractional-part-gaps} We have
\[
 x_{i+1}-x_i\le\frac5q\qquad(i\in[q]_{\Z}).
\]
\item\label{c:fractional-part-count} We have
\[
 \#\{j\in[0,q-1]_{\Z}:\{w+j\sigma\}\in[0,2\eta]\cup[1-2\eta,1)\}
 \le4\eta q+10.
\]
\end{enumerate}
\end{lemma}
\begin{proof}
\prf{Claim~\ref{c:fractional-part-gaps}} Let us compare the fractional parts with a rational grid.
For $j\in[0,q-1]_{\Z}$, put $t_j=jp-q\floor{w+j\sigma}$.
Then,
\[
 \left|\{w+j\sigma\}-\left(w+\frac{t_j}{q}\right)\right|
 =j\left|\sigma-\frac pq\right|<\frac2q.
\]
The integers $t_j$ represent distinct residue classes modulo $q$ because
$t_j\equiv jp\pmod q$ and $\gcd(p,q)=1$.
Thus, each $t\in\Z$ has the form $t=t_j+nq$ for some
$j\in[0,q-1]_{\Z}$ and $n\in\Z$.
Suppose $x_{i+1}-x_i>\frac5q$.
The interval $(x_i+\frac2q,x_{i+1}-\frac2q)$ has length greater than $\frac1q$.
The interval contains $w+\frac tq$ for some $t\in\Z$.
Let $t=t_j+nq$.
The number $\{w+j\sigma\}+n$ belongs to $(x_i,x_{i+1})$,
contradicting the choice of $x_i$ and $x_{i+1}$.

\prf{Claim~\ref{c:fractional-part-count}} Let us count the grid points near the two endpoints.
If
$\{w+j\sigma\}\in[0,2\eta]\cup[1-2\eta,1)$, then
\[
 w+\frac{t_j}{q}\in
 \left[-\frac2q,2\eta+\frac2q\right]
 \cup\left[1-2\eta-\frac2q,1+\frac2q\right].
\]
Each interval has length $2\eta+\frac4q$ and contains at most $2\eta q+5$
numbers of the form $w+\frac tq$, $t\in\Z$.
Hence, at most $4\eta q+10$ indices satisfy the condition in Claim~\ref{c:fractional-part-count}.
\end{proof}

\Needspace{20\baselineskip}
\begin{lemma}[Phase alignment and prefix discrepancies]\label{lem:phase-prefix-discrepancy}
Let $p_0\in\Z_{>0}$, $q_0\in\Z_{>0}$, $L\in\Z_{>0}$, $0<\eta<\frac14$, $\sigma_1\in[1,\infty)$, and $\sigma_2\in[1,\infty)$. Assume
\[
 \gcd(p_0,q_0)=1,\qquad
 q_0>\frac{20}{\eta},\qquad q_0\le\eta L,
\]
\[
 \left|\sigma_a-\frac{p_0}{q_0}\right|<\frac{2}{q_0^2}
 \quad(a\in[2]_{\Z}),\qquad
 |\sigma_1-\sigma_2|<\frac{\eta}{2L}.
\]
Let $u_1\in[0,1)$, $u_2\in[0,1)$, and $b\in\Z_{\ge0}$. For $j\in[0,q_0-1]_{\Z}$, define
\[
 v_j=\{u_2+(b+j)\sigma_2\},\qquad
 J=\{j\in[0,q_0-1]_{\Z}:|u_1-v_j|<\eta\}.
\]
The following statements hold.

\begin{enumerate}
\item\label{c:phase-selection} The set $J$ is nonempty. Define the uniquely specified index and phase
\[
j_* =\min J,\qquad v=v_{j_*}.
\]
Then, $0\le j_*<q_0$ and $|u_1-v|<\eta$.

\item\label{c:phase-prefix-mass} Use $v$ from Claim~\ref{c:phase-selection}. Define
\[
d_i=\lfloor u_1+i\sigma_1\rfloor-\lfloor v+i\sigma_2\rfloor
\quad(i\in[0,L]_{\Z}).
\]
We have $d_0=0$. For $i\in[L]_{\Z}$,
\[
d_i\in\{-1,0,1\},\qquad
d_i\ne0\Longrightarrow
\{u_1+i\sigma_1\}\in[0,2\eta]\cup[1-2\eta,1).
\]
Moreover,
\[
\sum_{i=1}^L|d_i|
\le\left(4\eta+\frac{10}{q_0}\right)L+q_0
\le\frac{11}{2}\eta L.
\]

\item\label{c:phase-run-variation} Define, for $i\in[L]_{\Z}$,
\[
\alpha_i=\lfloor u_1+i\sigma_1\rfloor-\lfloor u_1+(i-1)\sigma_1\rfloor,
\qquad
\beta_i=\lfloor v+i\sigma_2\rfloor-\lfloor v+(i-1)\sigma_2\rfloor.
\]
Then,
\[
\alpha_i-\beta_i=d_i-d_{i-1},\qquad
\sum_{i=1}^L|\alpha_i-\beta_i|
\le2\sum_{i=1}^L|d_i|-|d_L|
\le11\eta L.
\]
\end{enumerate}
\end{lemma}
\begin{proof}
\prf{Claim~\ref{c:phase-selection}} Let us prove that $J$ is nonempty. Let $x_1\le\cdots\le x_{q_0}$ be the ordered values of $v_0,\ldots,v_{q_0-1}$. Let $x_{q_0+1}=x_1+1$.
Claim~\ref{c:fractional-part-gaps} of Lemma~\ref{lem:fractional-part-estimates} with phase $u_2+b\sigma_2$ gives
\[
 x_{a+1}-x_a\le\frac5{q_0}<\eta
 \quad(a\in[q_0]_{\Z}).
\]
If $x_a\le u_1\le x_{a+1}$ for some $a<q_0$, then $|u_1-x_a|\le\frac5{q_0}<\eta$. If $u_1\le x_1$, use
\[
 x_1-u_1\le x_1\le x_1+1-x_{q_0}\le\frac5{q_0}.
\]
If $u_1\ge x_{q_0}$, use
\[
 u_1-x_{q_0}<1-x_{q_0}\le x_1+1-x_{q_0}\le\frac5{q_0}.
\]
Every $x_i$ equals $v_j$ for some $j\in[0,q_0-1]_{\Z}$. Hence, $J\ne\varnothing$. Since $j_*=\min J\in J$, we have $0\le j_*<q_0$ and $|u_1-v_{j_*}|<\eta$.

\prf{Claim~\ref{c:phase-prefix-mass}} Let us locate every nonzero discrepancy. Since $u_1\in[0,1)$ and $v\in[0,1)$, we have $d_0=0$. For $i\in[L]_{\Z}$, put
\[
 a_i=u_1+i\sigma_1,\qquad b_i=v+i\sigma_2.
\]
Claim~\ref{c:phase-selection} gives $|u_1-v|<\eta$.
Together with $|\sigma_1-\sigma_2|<\frac{\eta}{2L}$ and $i\le L$, this gives
\[
 |a_i-b_i|
 \le|u_1-v|+i|\sigma_1-\sigma_2|
 <\eta+\frac{\eta}{2}=\frac{3\eta}{2}<1.
\]
Hence, the integer $d_i=\lfloor a_i\rfloor-\lfloor b_i\rfloor$ belongs to $\{-1,0,1\}$. If $d_i=1$, then $b_i<\lfloor a_i\rfloor$. Thus,
\[
 \{a_i\}=a_i-\lfloor a_i\rfloor<a_i-b_i<\frac{3\eta}{2}<2\eta.
\]
If $d_i=-1$, then $b_i\ge\lfloor a_i\rfloor+1$. Thus,
\[
 \{a_i\}\ge1-(b_i-a_i)>1-\frac{3\eta}{2}>1-2\eta.
\]
Thus, $d_i\ne0$ implies $\{u_1+i\sigma_1\}\in[0,2\eta]\cup[1-2\eta,1)$.

Let us count these indices. Write
\[
 H=\left\lfloor\frac{L}{q_0}\right\rfloor,\qquad
 s=L-Hq_0\in[0,q_0-1]_{\Z}.
\]
Partition $[L]_{\Z}$ into
\[
 B_z=[zq_0+1,(z+1)q_0]_{\Z}
 \quad(z\in[0,H-1]_{\Z})
\]
and the final interval $[Hq_0+1,L]_{\Z}$, which has $s$ elements. For $i=zq_0+1+j$, $j\in[0,q_0-1]_{\Z}$, let
\[
 w_z=u_1+(zq_0+1)\sigma_1.
\]
Then, $\{u_1+i\sigma_1\}=\{w_z+j\sigma_1\}$. Claim~\ref{c:fractional-part-count} of Lemma~\ref{lem:fractional-part-estimates} bounds the number of nonzero $d_i$ in $B_z$ by $4\eta q_0+10$. Since $|d_i|$ equals $1$ precisely at the nonzero discrepancies,
\[
 \sum_{i=1}^L|d_i|
 \le H(4\eta q_0+10)+s
 \le\left(4\eta+\frac{10}{q_0}\right)L+q_0.
\]
The inequalities $\frac{10}{q_0}<\frac\eta2$ and $q_0\le\eta L$ give the bound $\frac{11}{2}\eta L$. The estimate $\sum_{i=1}^L|d_i|\le H(4\eta q_0+10)+s$ also holds when $s=0$ and $[Hq_0+1,L]_{\Z}=\varnothing$.

\prf{Claim~\ref{c:phase-run-variation}} Let us express $\alpha_i-\beta_i$ as $d_i-d_{i-1}$ for $i\in[L]_{\Z}$. Expanding the definitions gives $\alpha_i-\beta_i=d_i-d_{i-1}$. Since $d_0=0$, the triangle inequality gives
\[
 \begin{aligned}
 \sum_{i=1}^L|\alpha_i-\beta_i|
 &\le\sum_{i=1}^L|d_i|+\sum_{i=1}^L|d_{i-1}|\\
 &=2\sum_{i=1}^L|d_i|-|d_L|\\
 &\le11\eta L.
 \end{aligned}
\]
The inequality $2\sum_{i=1}^L|d_i|-|d_L|\le11\eta L$ follows from Claim~\ref{c:phase-prefix-mass}.
\end{proof}

\begin{lemma}[Mechanical run-edits]\label{lem:mechanical-edits}
Let $m\in\Z_{\ge4}$ and $M=m-2$.
Let $r\in(1,M)$ be irrational and let $0<\eta<\frac14$.
Then, there are constants $\delta_m(r,\eta)>0$ and $B_m(r,\eta)>0$ depending only on $m$, $r$, and $\eta$ with the following property.
For each $(M,r,\eta,\delta_m(r,\eta),B_m(r,\eta))$-admissible tuple
$(l,T,\sigma_1,\sigma_2,u_1,u_2,(A_1,\ldots,A_l))$ and $U_0=A_1\ccn\cdots\ccn A_l$,
there are $n\in\Z_{\ge0}$, $\zeta\in\Z^2$, and a run-edit list $\Omega=(\omega_1,\ldots,\omega_n)$ with the following properties.
\begin{enumerate}
\item\label{c:mechanical-target} The list $\Omega$ starts at $U_0$ and satisfies
\[
 \operatorname{RunEResult}(\Omega,U_0)
 =\operatorname{Trans}\bigl(\operatorname{RunFactor}_{1,T}(\operatorname{MechRun}(\sigma_2,u_2)),\zeta\bigr).
\]
\item\label{c:mechanical-counts} $\Omega$ consists of $K=T-l$ insertions of simple runs,
$I$ vertical-edge insertions, and $D$ vertical-edge deletions. Their numbers satisfy
\begin{equation}\label{eq:mechanical-bound}
 n=K+I+D,\qquad I,D\in\Z_{\ge0},\qquad I+D\le12\eta l+B_m(r,\eta).
\end{equation}
\item\label{c:mechanical-lengths} For $U_j=\operatorname{RunEResult}((\omega_1,\ldots,\omega_j),U_0)$ $(j\in[0,n]_{\Z})$,
every simple run $R\subseteq U_j$ satisfies $|R|_N\in[M]_{\Z}$.
\end{enumerate}
\end{lemma}
\begin{proof}
\prf{Claim~\ref{c:mechanical-target}} Let us choose $p_0$, $q_0$, $L$, $\delta_m(r,\eta)$, and $B_m(r,\eta)$.
Put $k=\floor r$.
Choose relatively prime $p_0\in\Z_{>0}$ and $q_0\in\Z_{>0}$ such that
\begin{equation}\label{eq:mechanical-approximation}
 q_0>\frac{20}{\eta},\qquad |r-\frac{p_0}{q_0}|<\frac1{q_0^2}.
\end{equation}
To obtain such a fraction, partition $[0,1)$ into $Q$ intervals of length $\frac1Q$.
Two of the $Q+1$ numbers $\{jr\}$, $j\in[0,Q]_{\Z}$, lie in the same interval.
Subtracting their indices gives integers $p$ and $q$ such that
\[
 1\le q\le Q,\qquad |qr-p|<\frac1Q,\qquad
 \left|r-\frac pq\right|<\frac1{qQ}\le\frac1{q^2}.
\]
Let $\frac{p'}{q'}$ be $\frac pq$ in lowest terms. Then, $1\le q'\le q$ and $\left|r-\frac{p'}{q'}\right|<\frac1{q^2}\le\frac1{q'^2}$.
These reduced denominators are unbounded as $Q\to\infty$: otherwise, only finitely many
reduced fractions in $[r-1,r+1]$ could occur, and their positive distances from the irrational number $r$
could not tend to zero.
Thus, choose a reduced fraction with $q_0>\frac{20}{\eta}$ satisfying~\eqref{eq:mechanical-approximation}.
Choose $L\in\Z_{>0}$ such that $q_0\le\eta L$ and let
\begin{equation}\label{eq:mechanical-constants}
 \delta_m(r,\eta)=\frac12\min\left\{\frac1{q_0^2},\frac{\eta}{4L},r-k,k+1-r\right\},
 \qquad B_m(r,\eta)=(M-1)L+1.
\end{equation}
These constants depend only on $r$, $\eta$, and $M$.
Condition~\ref{hyp:mechanical-proximity} of Definition~\ref{def:admissible-mechanical} gives $k<\sigma_i<k+1$ for $i\in[2]_{\Z}$.
Thus, every simple run of either mechanical sequence has vertical length $k$ or $k+1$. Both lengths belong to $[M]_{\Z}$.

Let us compare the two slopes to the same rational approximation.
For $\sigma\in\{\sigma_1,\sigma_2\}$, equations~\eqref{eq:mechanical-approximation}--\eqref{eq:mechanical-constants} give
\[
 |\sigma-\frac{p_0}{q_0}|<\frac2{q_0^2}.
\]
Thus, Lemma~\ref{lem:fractional-part-estimates} applies with $p=p_0$ and $q=q_0$.

Let us partition the run into factors of length $L$ and one final factor.
Let $H=\floor{\frac{l-1}{L}}$.
Let us construct the part of $\Omega$ acting on the first $HL$ simple runs of $U_0$.
Put $c_0=0$.
For $z\in[0,H-1]_{\Z}$, suppose the run-edits on the first $zL$ simple runs have
produced the first $c_z$ simple runs of $\operatorname{MechRun}(\sigma_2,u_2)$.

Let us choose the target intercept for the factor indexed by $z\in[0,H-1]_{\Z}$.
Let $u=\{u_1+zL\sigma_1\}$.
Apply Claim~\ref{c:phase-selection} of Lemma~\ref{lem:phase-prefix-discrepancy}
with source phase $u$ and offset $b=c_z$.
Let $j_z=\min J$ and $v=\{u_2+(c_z+j_z)\sigma_2\}$ as in Claim~\ref{c:phase-selection} of Lemma~\ref{lem:phase-prefix-discrepancy}.
Then, $0\le j_z<q_0$ and $|u-v|<\eta$.
If $j_z>0$, insert
$\operatorname{RunFactor}_{c_z+1,j_z}(\operatorname{MechRun}(\sigma_2,u_2))$
immediately before the copy of $\operatorname{RunFactor}_{zL+1,L}(U_0)$.

Let us bound the number of vertical run-edits on the copy of $\operatorname{RunFactor}_{zL+1,L}(U_0)$.
Define
\[
 X=\operatorname{RunFactor}_{1,L}(\operatorname{MechRun}(\sigma_1,u)),\qquad
 Y=\operatorname{RunFactor}_{1,L}(\operatorname{MechRun}(\sigma_2,v)).
\]
Claim~\ref{c:phase-run-variation} of Lemma~\ref{lem:phase-prefix-discrepancy} gives
\[
 \sum_{i=1}^L
 \left|\,|\operatorname{RunFactor}_{i,1}(X)|_N
       -|\operatorname{RunFactor}_{i,1}(Y)|_N\,\right|\le11\eta L.
\]
By Lemma~\ref{lem:vertical-length-comparison}, there is a run-edit list $\Omega_z=(\omega_{z,1},\ldots,\omega_{z,n_z})$ satisfying $\operatorname{RunEResult}(\Omega_z,X)=Y$ and consisting of at most $11\eta L$
vertical run-edits.
Every simple-run vertical length produced by $\Omega_z$ belongs to $[M]_{\Z}$.
Let $c_{z+1}=c_z+j_z+L$.
For some $\zeta_z\in\Z^2$, the resulting first $c_{z+1}$ simple runs form
\[
 \operatorname{Trans}\bigl(\operatorname{RunFactor}_{1,c_{z+1}}(\operatorname{MechRun}(\sigma_2,u_2)),\zeta_z\bigr).
\]

Let us construct the run-edits on $A_{HL+1},\ldots,A_l$.
We have $l-HL\in[L]_{\Z}$.
For $i\in[l-HL]_{\Z}$, change the length of the copy of $A_{HL+i}$ to
$|\operatorname{RunFactor}_{c_H+i,1}(\operatorname{MechRun}(\sigma_2,u_2))|_N$.
Both lengths belong to $[M]_{\Z}$, including the length of $A_l$ by Condition~\ref{hyp:mechanical-source} of Definition~\ref{def:admissible-mechanical}.
Lemma~\ref{lem:vertical-length-comparison} gives at most $(M-1)L$ vertical run-edits
and keeps every length in $[M]_{\Z}$.
Now,
\[
 c_H+l-HL=l+\sum_{z=0}^{H-1}j_z\le l+Hq_0\le(1+\eta)l<T.
\]
Append
\[
 \operatorname{RunFactor}_{c_H+l-HL+1,\,T-c_H-l+HL}
       (\operatorname{MechRun}(\sigma_2,u_2)).
\]
Let $Q=\operatorname{RunEResult}(\Omega,U_0)$ and $\zeta=\operatorname{Start}_{\mathrm{pt}}(Q)$.
Then,
\[
 Q=\operatorname{Trans}\bigl(\operatorname{RunFactor}_{1,T}(\operatorname{MechRun}(\sigma_2,u_2)),\zeta\bigr).
\]

\prf{Claim~\ref{c:mechanical-counts}} Let us count the simple-run insertions and the vertical run-edits.
The construction in Claim~\ref{c:mechanical-target} uses exactly $K=T-l$ simple-run insertions. If $I$ and $D$ count its vertical-edge insertions and deletions, respectively, then $n=K+I+D$. The bound in~\eqref{eq:mechanical-bound} follows from
\[
 I+D\le11\eta HL+(M-1)L\le12\eta l+B_m(r,\eta).
\]
\prf{Claim~\ref{c:mechanical-lengths}} Every inserted simple run has length $k$ or $k+1$. Both belong to $[M]_{\Z}$.
Lemma~\ref{lem:vertical-length-comparison} keeps every intermediate vertical length in $[M]_{\Z}$ at each application in the construction of $\Omega$.
\end{proof}

\begin{lemma}\label{lem:approaching-sequence}
Let $m\in\Z_{\ge3}$.
For every $r\in I_m$, there is a sequence $((a_\nu,b_\nu))_{\nu\ge1}$ in $\cA_m$ satisfying
\begin{equation}\label{eq:approaching-sequence}
 \lim_{\nu\to\infty}a_\nu=\infty,\qquad
 \lim_{\nu\to\infty}\frac{b_\nu}{a_\nu}=r.
\end{equation}
\end{lemma}
\begin{proof}
Let us choose pairs in $\cA_m$ whose slopes tend to $r$.
Choose distinct reduced fractions $\frac{b_\nu}{a_\nu}\in I_m$ tending to $r$.
For each $A>0$, the bounded interval $I_m$ contains only finitely many reduced
fractions with denominator at most $A$. Thus, $a_\nu\to\infty$.
Since $\frac{b_\nu}{a_\nu}\in I_m$, we have
$a_\nu^2+b_\nu^2-ma_\nu b_\nu\le0$.
Equation~\eqref{eq:nonempty-domain} gives $(a_\nu,b_\nu)\in\cA_m$.
\end{proof}

\begin{theorem}\label{thm:irrational-large-m}
Let $m\in\Z_{\ge4}$ and $M=m-2$.
For every irrational $r\in(1,M)$, there is $\Upsilon_m(r)\in\R$ such that
\[
 \lim_{\substack{a\to\infty,\ (a,b)\in\cA_m\\\frac{b}{a}\to r}}
 \frac{\log E_m(a,b)}a=\Upsilon_m(r).
\]
\end{theorem}
\begin{proof}
Let us choose the comparison parameters for Lemma~\ref{lem:mechanical-edits} and Claim~\ref{c:composition-bound} of Proposition~\ref{prop:composition}.
Let
\[
 C=32(m+1),\qquad H=2C+16,
 \qquad 0<\eta<\frac1{2H},\qquad d=1-H\eta.
\]
Use $\delta_m(r,\eta)$ and $B_m(r,\eta)$ from Lemma~\ref{lem:mechanical-edits}.
Choose relatively prime $p\in\Z_{>0}$ and $q\in\Z_{>0}$ such that $q\ge2$,
$\alpha=\frac{p}{q}\in(1,M)$, and $|\alpha-r|<\delta_m(r,\eta)$.
We fix $\alpha$, $p$, and $q$ while $a\to\infty$, $(a,b)\in\cA_m$, and $\frac ba\to r$.

Let us verify~\eqref{eq:layer-budget} for Claim~\ref{c:composition-bound} of Proposition~\ref{prop:composition}.
For a run $Z$ of $\nu$ simple runs and sufficiently large $\nu$, let
\[
 P=\floor{C\eta\nu},\qquad T=\nu-2P-2,\qquad
 V=\operatorname{RunFactor}_{P+2,T}(Z).
\]
We have $P>0$ and $T>0$. Define
\[
\begin{aligned}
 W_{\mathrm L}&=\operatorname{RunFactor}_{1,P}(Z),&
 W_{\mathrm R}&=\operatorname{RunFactor}_{\nu-P+1,P}(Z),\\
 u&=|\operatorname{RunFactor}_{P+1,1}(Z)|_N,&
 v&=|\operatorname{RunFactor}_{\nu-P,1}(Z)|_N.
\end{aligned}
\]
Equation~\eqref{eq:target-shape} gives $Z=W_{\mathrm L}\ccn R_u\ccn V\ccn R_v\ccn W_{\mathrm R}$.
Let $l=d\nu+O(1)$. Then,
\[
 K=T-l,\qquad \lim_{\nu\to\infty}\frac K\nu=16\eta,\qquad
 T\ge(1+\eta)l+B_m(r,\eta)
\]
when $\nu$ is sufficiently large. If $I\in\Z_{\ge0}$ and $D\in\Z_{\ge0}$ satisfy $I+D\le13\eta\nu$, then eventually
\[
 K+I+(m+1)D\le(17+13(m+1))\eta \nu\le P.
\]
The inequality $(17+13(m+1))\eta\nu\le P$ uses $P=\floor{C\eta\nu}$ and
$(C-17-13(m+1))\eta\nu\ge1$ for sufficiently large $\nu$, since
$C=32(m+1)>17+13(m+1)$. Equation~\eqref{eq:entropy-loss} gives
\[
 \log \operatorname{Fib}_m(l,D)\le\nu\psi_m(13\eta).
\]

Let us verify admissibility in the two applications below.
Equation~\eqref{eq:ambient} gives $\Gamma(a,b)\to r$.
Thus, $1<\Gamma(a,b)<M$ and $|\Gamma(a,b)-r|<\delta_m(r,\eta)$ for all sufficiently large $a$.
The choice of $\alpha$ and the bound $T\ge(1+\eta)l+B_m(r,\eta)$ verify condition~\ref{hyp:mechanical-proximity} of Definition~\ref{def:admissible-mechanical} in both applications.

First, let us obtain a lower bound for $\frac1a\log E_m(a,b)$.
Let $(a,b)\in\cA_m$, $\nu=a+1$, and $j=\floor{\frac{d\nu}{q}}$.
Apply Lemma~\ref{lem:mechanical-edits} with
\[
\begin{gathered}
 U=W_{p,q}^j,\quad l=jq,\quad Z=\mathcal U(a,b),\quad
 V=\operatorname{RunFactor}_{P+2,T}(\mathcal U(a,b)),\\
 (\sigma_1,\sigma_2,u_1,u_2)
 =\bigl(\alpha,\Gamma(a,b),0,\{(P+1)\Gamma(a,b)\}\bigr),\\
 A_i=R_{w_i(p,q)}\qquad(i\in[l]_{\Z}).
\end{gathered}
\]
Definition~\ref{def:mechanical-paths} verifies condition~\ref{hyp:mechanical-source} of Definition~\ref{def:admissible-mechanical} for these source runs.
Lemma~\ref{lem:mechanical-edits} gives $I+D\le12\eta l+B_m(r,\eta)\le13\eta\nu$
for all sufficiently large $a$.
The last run $\operatorname{RunFactor}_{a+1,1}(\mathcal U(a,b))$
belongs to $W_{\mathrm R}$ and does not occur in $V$.
Claim~\ref{c:composition-bound} of Proposition~\ref{prop:composition}
and Lemma~\ref{lem:word-reindexing} apply with $V_*=\mathcal U(a,b)$.
Equations~\eqref{eq:dictionary}, \eqref{eq:composition}, and~\eqref{eq:entropy-loss} give
\[
 \frac{\log E_m(a,b)}a
 \ge\frac{j}{a}\frac{\log C_0(W_{p,q}^j)}j
 -\frac{\nu}{a}\psi_m(13\eta)-\frac{\log m}{a}.
\]
Since $\frac ja\to\frac dq$ and $\frac\nu a\to1$, Claim~\ref{c:periodic-limit} of Proposition~\ref{prop:periodic} and~\eqref{eq:rational-period} give
\begin{equation}\label{eq:irrational-lower}
 \liminf_{\substack{a\to\infty,\ (a,b)\in\cA_m\\\frac{b}{a}\to r}}
 \frac{\log E_m(a,b)}a
 \ge d\Upsilon_m(\alpha)-\psi_m(13\eta).
\end{equation}

Second, let us obtain an upper bound for $\frac1a\log E_m(a,b)$.
Let $(a,b)\in\cA_m$, $l=a+1$, $N=\ceil{\frac{l}{dq}}$, and $\nu=Nq$.
Apply Lemma~\ref{lem:mechanical-edits} with
\[
\begin{gathered}
 U=\mathcal U(a,b),\quad Z=W_{p,q}^N,\quad
 V=\operatorname{RunFactor}_{P+2,T}(W_{p,q}^N),\\
 (\sigma_1,\sigma_2,u_1,u_2)
 =\bigl(\Gamma(a,b),\alpha,0,\{(P+1)\alpha\}\bigr),\\
 A_i=R_{\Lambda_i(a,b)}\quad(i\in[a]_{\Z}),\qquad A_l=R_1.
\end{gathered}
\]
Definition~\ref{def:frame} verifies condition~\ref{hyp:mechanical-source} of Definition~\ref{def:admissible-mechanical}, including $|A_l|_N=1\in[M]_{\Z}$.
Lemma~\ref{lem:mechanical-edits} gives $I+D\le12\eta l+B_m(r,\eta)\le13\eta\nu$
for all sufficiently large $a$.
Claim~\ref{c:composition-bound} of Proposition~\ref{prop:composition}
and Lemma~\ref{lem:word-reindexing} apply with $V_*=W_{p,q}^N$.
Equations~\eqref{eq:composition}, \eqref{eq:dictionary}, and~\eqref{eq:entropy-loss} give
\[
 \frac{\log E_m(a,b)}a
 \le\frac{N}{a}\frac{\log C_0(W_{p,q}^N)}N
 +\frac{\nu}{a}\psi_m(13\eta)-\frac{\log m}{a}.
\]
Here, $\frac Na\to\frac1{dq}$ and $\frac\nu a\to\frac1d$. By Claim~\ref{c:periodic-limit} of Proposition~\ref{prop:periodic} and~\eqref{eq:rational-period},
\begin{equation}\label{eq:irrational-upper}
 \limsup_{\substack{a\to\infty,\ (a,b)\in\cA_m\\\frac{b}{a}\to r}}
 \frac{\log E_m(a,b)}a
 \le\frac{\Upsilon_m(\alpha)+\psi_m(13\eta)}d.
\end{equation}


Third, let us compare~\eqref{eq:irrational-lower} and~\eqref{eq:irrational-upper}. Lemma~\ref{lem:word-bounds} gives
$\Upsilon_m(\alpha)\le m(M+1)\log2$.
The interval in~\eqref{eq:irrational-lower}--\eqref{eq:irrational-upper} has width at most
\[
 m(M+1)\log2\left(\frac1d-d\right)
      +\psi_m(13\eta)\left(1+\frac1d\right).
\]
As $\eta\downarrow0$, we have $d=1-H\eta\to1$ and
\[
 \frac1d-d\longrightarrow0,\qquad
 \psi_m(13\eta)\longrightarrow0.
\]
The limit $\psi_m(13\eta)\to0$ as $\eta\downarrow0$ follows from $x\log x\to0$ as $x\downarrow0$ and Definition~\ref{def:entropy-loss-function}.

Let us prove that every sequence has the same limit.
For each $i\in[2]_{\Z}$, let $((a_{i,n},b_{i,n}))_{n\ge1}$ be a sequence in $\cA_m$ satisfying
\[
 \lim_{n\to\infty}a_{i,n}=\infty,\qquad
 \lim_{n\to\infty}\frac{b_{i,n}}{a_{i,n}}=r.
\]
Define
\[
 L_i=\liminf_{n\to\infty}\frac{\log E_m(a_{i,n},b_{i,n})}{a_{i,n}},\qquad
 U_i=\limsup_{n\to\infty}\frac{\log E_m(a_{i,n},b_{i,n})}{a_{i,n}}.
\]
Lemma~\ref{lem:word-bounds} and~\eqref{eq:dictionary} make these four numbers finite.

Let us prove $U_i\le L_j$ for $i,j\in[2]_{\Z}$.
For each fixed $\eta$, choose $\alpha$, $\delta_m(r,\eta)$, and $B_m(r,\eta)$ before taking either sequence limit.
Equations~\eqref{eq:irrational-lower} and~\eqref{eq:irrational-upper} give, for $i,j\in[2]_{\Z}$,
\[
 U_i-L_j\le m(M+1)\log2\left(\frac1d-d\right)
      +\psi_m(13\eta)\left(1+\frac1d\right).
\]
Letting $\eta\downarrow0$ gives $U_i\le L_j$.
Taking $i=j$ and using $L_i\le U_i$ gives $L_i=U_i$.
Taking $(i,j)=(1,2)$ and $(2,1)$ gives $L_1=U_1=L_2=U_2$.
Lemma~\ref{lem:approaching-sequence} supplies at least one such sequence.
Define $\Upsilon_m(r)$ to be its limit. The equality $L_1=U_1=L_2=U_2$ makes this definition independent of the chosen sequence.
\end{proof}

\subsection{Continuity of the limit}
\begin{lemma}\label{lem:continuity}
Let $m\in\Z_{\ge3}$.
Let $h_1\in\R$ and $h_2\in\R$ satisfy $h_1<h_2$. Let $J=[h_1,h_2]$.
Let $F:J\to\R$. Assume the following.
\begin{hypotheses}
\item\label{hyp:all-sequence-limit} For each $r\in J$ and every sequence $((a_n,b_n))_{n\ge1}$ in $\cA_m$ satisfying
\[
 \lim_{n\to\infty}a_n=\infty,\qquad
 \lim_{n\to\infty}\frac{b_n}{a_n}=r,
\]
we have
\[
 \lim_{n\to\infty}\frac{\log E_m(a_n,b_n)}{a_n}=F(r).
\]
\item\label{hyp:approaching-sequences} For each $r\in J$, there exists a sequence $((a_n,b_n))_{n\ge1}$ in $\cA_m$ such that
\[
 \lim_{n\to\infty}a_n=\infty,\qquad
 \lim_{n\to\infty}\frac{b_n}{a_n}=r.
\]
\end{hypotheses}
Then, we have the following.

\begin{enumerate}
\item\label{c:continuity-neighborhood} For every $r\in J$ and $\varepsilon>0$, there exist $A>0$ and $\delta>0$ such that
\begin{equation}\label{eq:continuity-neighborhood}
 \left|\frac{\log E_m(a,b)}a-F(r)\right|<\varepsilon
 \quad\left((a,b)\in\cA_m,\ a>A,\ \left|\frac ba-r\right|<\delta\right).
\end{equation}
\item\label{c:continuity-function} The function $F$ is continuous on $J$.
\item\label{c:continuity-endpoint} At the endpoints of $J$, we have the following.
\begin{enumerate}[label=(\alph*),ref=\theenumi(\alph*)]
\item\label{c:continuity-endpoint-neighborhood} For $r\in\{h_1,h_2\}$, the estimate in Claim~\ref{c:continuity-neighborhood} holds when additionally $\frac ba\in J$.
\item\label{c:continuity-endpoint-function} We have
\[
 \lim_{s\downarrow h_1}F(s)=F(h_1),\qquad
 \lim_{s\uparrow h_2}F(s)=F(h_2)\qquad(s\in J).
\]
\end{enumerate}
\end{enumerate}
\end{lemma}
\begin{proof}
\prf{Claim~\ref{c:continuity-neighborhood}} Let us prove~\eqref{eq:continuity-neighborhood} by contradiction.
If~\eqref{eq:continuity-neighborhood} fails at $r$, there is $\varepsilon_0>0$
such that, for every positive integer $\nu$, we find $(a_\nu,b_\nu)\in\cA_m$ with
\[
 a_\nu>\nu,\qquad
 \left|\frac{b_\nu}{a_\nu}-r\right|<\frac1\nu,\qquad
 \left|\frac{\log E_m(a_\nu,b_\nu)}{a_\nu}-F(r)\right|\ge\varepsilon_0.
\]
We have $a_\nu\to\infty$ and $\frac{b_\nu}{a_\nu}\to r$.
Assumption~\ref{hyp:all-sequence-limit} contradicts
$|\frac{\log E_m(a_\nu,b_\nu)}{a_\nu}-F(r)|\ge\varepsilon_0$.

\prf{Claim~\ref{c:continuity-function}} Let us obtain continuity from~\eqref{eq:continuity-neighborhood}.
Let $r\in J$ and $\varepsilon>0$.
Claim~\ref{c:continuity-neighborhood} gives $\delta>0$ and $A>0$ such that
\[
 \left|\frac{\log E_m(a,b)}a-F(r)\right|<\frac\varepsilon2
 \quad\left((a,b)\in\cA_m,\ a>A,\ \left|\frac ba-r\right|<\delta\right).
\]

Let us prove $|F(s)-F(r)|<\varepsilon$ when $s\in J$ and $|s-r|<\frac\delta2$.
Fix such an $s$.
By assumption~\ref{hyp:approaching-sequences}, choose a sequence $((a_\nu,b_\nu))_{\nu\ge1}$ in $\cA_m$ such that
\[
 a_\nu\longrightarrow\infty,\qquad \frac{b_\nu}{a_\nu}\longrightarrow s
 \qquad(\nu\longrightarrow\infty).
\]
Choose $\nu_0\in\Z_{>0}$ such that
\[
 a_\nu>A,\qquad
 \left|\frac{b_\nu}{a_\nu}-s\right|<\frac\delta2
 \qquad(\nu\ge\nu_0).
\]
For $\nu\ge\nu_0$, the triangle inequality gives
\[
 \left|\frac{b_\nu}{a_\nu}-r\right|
 \le\left|\frac{b_\nu}{a_\nu}-s\right|+|s-r|<\delta.
\]
Hence,~\eqref{eq:continuity-neighborhood} applies with $\varepsilon$ replaced by $\frac\varepsilon2$ for every $\nu\ge\nu_0$.

Let us take the limit as $\nu\to\infty$.
Assumption~\ref{hyp:all-sequence-limit} and continuity of the absolute value give
\[
 |F(s)-F(r)|
 =\lim_{\nu\to\infty}\left|\frac{\log E_m(a_\nu,b_\nu)}{a_\nu}-F(r)\right|
 \le\frac\varepsilon2<\varepsilon.
\]

\prf{Claim~\ref{c:continuity-endpoint-neighborhood}}
Claim~\ref{c:continuity-neighborhood} applies to $r\in\{h_1,h_2\}$ and pairs satisfying $\frac ba\in J$.

\prf{Claim~\ref{c:continuity-endpoint-function}} Let us identify the one-sided endpoint limits.
The hypothesis $h_1<h_2$ makes both one-sided domains nonempty.
By Claim~\ref{c:continuity-function}, for every $\varepsilon>0$ there are $\delta_1>0$ and $\delta_2>0$ such that
\[
\begin{aligned}
 |F(s)-F(h_1)|&<\varepsilon
 &&(s\in J,\ 0<s-h_1<\delta_1),\\
 |F(s)-F(h_2)|&<\varepsilon
 &&(s\in J,\ 0<h_2-s<\delta_2).
\end{aligned}
\]
These are the right limit at $h_1$ and the left limit at $h_2$, respectively.
\end{proof}


